% ---------------- ABSTRACT ----------------
\chapter*{Abstract}
\addcontentsline{toc}{chapter}{Abstract}

Large language models produce fluent text that is sometimes confidently wrong. This
failure, commonly called \emph{hallucination}, is now the dominant practical obstacle
to deploying such models where correctness matters. A promising line of detectors
reads the model's own internal activations as it generates, and learns to tell
truthful continuations from fabricated ones without any external retrieval and without
sampling the model many times. These detectors are attractive because they add almost
no cost at inference time, but the published results are scattered across different
models, datasets, and hardware budgets, which makes it hard to know how much each
ingredient actually contributes.

This dissertation studies internal-state hallucination detection under a single,
reproducible protocol that runs entirely on free cloud hardware. We adopt the
label-free data-generation idea of the MIND framework, in which a model continues a
truncated encyclopedia article and the continuation is automatically labelled by
whether the model recovers the original entity. On top of the canonical last-token,
last-layer hidden state used by that framework, we add three inexpensive signals
computed from the same forward pass: how much the representation shifts from one layer
to the next, how much it varies across layers, and how uncertain the model is while
generating. We then compare this augmented detector against a wider set of literature
features and against four established baselines, together with two simple reference
methods, across as many as ten question-answering, summarisation, and dialogue
benchmarks.

We report only the results that our saved measurement files support, and we separate
\emph{in-distribution} detection (held-out encyclopedia text) from \emph{cross-task
transfer} (the question-answering and dialogue benchmarks), because the two regimes
behave very differently. On the fully measured six-billion-parameter model, the
richest feature configurations are competitive with the strongest baseline but do not
overtake it, and a single benchmark heavily influences the averages; on a second model
of comparable size an earlier extended run favoured one of the proposed feature
configurations. The honest summary is that the value of these internal-state features
is real but \emph{model dependent}, and that cross-task transfer remains the hard part.
The contribution of this work is therefore not a new state-of-the-art detector but a
careful, free-tier-reproducible comparison of internal-state features across several
open models, with an explicit account of what is measured, what is still pending, and
where the method's limits lie.

\vspace{0.4cm}
\noindent\textbf{Keywords:} hallucination detection, large language models, internal
representations, hidden states, probing classifiers, model calibration, reproducible
evaluation.
